\documentclass[11pt]{article}
\usepackage[margin=0.85in]{geometry}
\usepackage{booktabs}
\usepackage{array}
\usepackage{amsmath}
\usepackage{float}
\usepackage{xcolor}
\usepackage[hidelinks]{hyperref}
\usepackage[numbers,sort&compress]{natbib}

\definecolor{accent}{HTML}{2E7180}
\definecolor{light}{HTML}{EAF3F5}
\newcolumntype{L}[1]{>{\raggedright\arraybackslash}p{#1}}

\title{Week 35 Draft: Affinity-Model Architecture and\\
RFdiffusion--ProteinMPNN Binder Design}
\author{Yusheng Cai}
\date{August 24--30, 2026}

\begin{document}
\maketitle
\begin{center}
  \small Repository:
  \href{https://github.com/Yusheng-cai/BindingAffinityPredictor}%
       {\texttt{github.com/Yusheng-cai/BindingAffinityPredictor}}
\end{center}

\section*{Executive summary}

This report will address two connected but scientifically distinct workflows.
The first is \textbf{affinity prediction}: Boltz-2, Nesso-1, and FlashBind map
protein--small-molecule information to an affinity score, with different use
of sequence, MSA, pocket, and three-dimensional coordinates. The second is
\textbf{protein-binder design}: RFdiffusion proposes a binder backbone around
a target protein and ProteinMPNN assigns amino-acid sequences compatible with
that backbone \citep{watson2023rfdiffusion,dauparas2022proteinmpnn}.

RFdiffusion and ProteinMPNN are included because they expose a useful location
for adding physical information during design. They are not affinity models,
and their native scores are not binding free energies or probabilities of
experimental success. I will keep their analysis separate from the
small-molecule affinity benchmark and use the audited insulin-receptor example
as a concrete worked case.

\section{Questions for the week}

\begin{enumerate}
  \item What tensors and geometric representations connect the input of each
  affinity model to its affinity head?
  \item What exactly do RFdiffusion and ProteinMPNN receive and return?
  \item How do hotspot conditioning and stochastic sampling alter the proposed
  binder backbone?
  \item Where could a physics-aware interface score enter the binder-design
  workflow without being confused with ProteinMPNN likelihood?
\end{enumerate}

\clearpage
\section{Two different prediction problems}

\begin{table}[H]
\centering
\footnotesize
\begin{tabular}{L{0.19\textwidth}L{0.23\textwidth}L{0.24\textwidth}L{0.21\textwidth}}
\toprule
Task & Representative input & Native output & Appropriate evaluation \\
\midrule
Small-molecule affinity prediction & Protein sequence or structure and ligand chemistry/pose & Affinity score, and for Boltz-2 a cofolded complex & Within-target affinity correlation; protein and ligand structure recovery when references exist \\
RFdiffusion binder generation & Target coordinates, contig, optional hotspot residues, checkpoint, and random seed & Target plus a generated binder backbone and reverse-diffusion trajectories & Backbone geometry, target-relative pose, contacts, clashes, diversity, and downstream validation \\
ProteinMPNN sequence design & Fixed target/binder backbone, design mask, checkpoint, temperature, and seed & Amino-acid sequences, per-residue probabilities, and sequence-model scores & Sequence diversity, structural compatibility, independent complex prediction, physical interface checks, and experiment \\
\bottomrule
\end{tabular}
\caption{The native outputs answer different questions. In particular,
ProteinMPNN likelihood is not an affinity, and RFdiffusion confidence is not a
binding measurement.}
\end{table}

\section{RFdiffusion: target-conditioned backbone generation}

\subsection{Input and output}

In the binder-design setting, RFdiffusion receives a three-dimensional target
structure and a contig specification that identifies fixed target residues and
the length or length range of the new chain. Optional hotspot residues nominate
the target region where an interface should be formed. The checkpoint and
random seed are also scientifically important inputs because different
checkpoints encode different training tasks and different seeds generate
different samples.

Only the binder is diffused in the target-conditioned protocol. The target
coordinates remain fixed and provide conditioning throughout reverse
diffusion. A simplified input--output statement is
\[
  (\text{target coordinates},\ \text{contig},\ \text{hotspots},\ \text{seed})
  \longrightarrow
  (\text{target} + \text{binder backbone},\ \text{trajectory}).
\]
The generated chain is initially represented as a backbone proposal; it is not
a chemically complete protein and it has no designed amino-acid sequence.

\subsection{Why equivariance matters}

RFdiffusion represents protein geometry using residue-level positions and
orientations. If every input coordinate is transformed by the same rigid
motion,
\[
  \mathbf{x}_i' = Q\mathbf{x}_i + \mathbf{a},
\]
the denoised output should transform by the same $Q$ and $\mathbf{a}$. This
SE(3)-equivariance prevents the model from attaching meaning to the arbitrary
orientation of a protein in a coordinate file. It also explains why geometric
guidance should depend on invariant distances, angles, contacts, or energies,
rather than laboratory-frame Cartesian coordinates.

\subsection{What hotspot conditioning establishes}

Hotspots are conditioning information, not residues that RFdiffusion discovers
as experimental binding sites. In the existing matched insulin-receptor
ablation, both arms used the same complex checkpoint, target, contig, ten
seeds, and sampling settings; only the hotspot tensor was removed. The guided
backbone had a smaller minimum C$\alpha$ distance to A59, A83, or A91 in 9 of
10 matched seeds. The median minimum distance was 6.38~\AA{} with hotspots and
9.72~\AA{} without them.

This demonstrates that hotspot information steers backbone placement. It does
not demonstrate affinity, folding, specificity, or validated pocket discovery.
The ten guided seeds also produced diverse backbones: pairwise intrinsic trace
RMSD had a median of 13.36~\AA, and target-aligned binder-pose Chamfer distance
had a median of 5.08~\AA. One RFdiffusion output is therefore one stochastic
proposal, not a unique solution.

\section{ProteinMPNN: sequence design on a fixed backbone}

ProteinMPNN addresses the next question: given fixed backbone coordinates,
which amino-acid identities are geometrically compatible with them
\citep{dauparas2022proteinmpnn}? It constructs a residue graph from backbone
geometry, passes messages using invariant geometric features, and samples a
sequence autoregressively. Schematically,
\[
  p(\mathbf{s}\mid X)
  = \prod_{k=1}^{L} p\!\left(s_{\pi_k}\mid
  s_{\pi_1},\ldots,s_{\pi_{k-1}},X\right),
\]
where $X$ is the fixed backbone, $\mathbf{s}$ is the sequence, and $\pi$ is a
decoding order. Target residues can be fixed while only the generated binder
chain is designed.

The model returns amino-acid probabilities and a sequence-model score. A low
negative log-probability means that ProteinMPNN considers the sequence
compatible with the supplied backbone under its learned distribution. It does
not mean low binding free energy, high affinity, or high experimental success.

For the fixed 90-residue RFdiffusion seed-42 backbone, the audited run generated
100 unique ProteinMPNN sequences at temperature 0.1. Median pairwise identity
was 66.7\%, median position entropy was 0.648 nats, and 17 of 90 positions were
invariant in this finite sample. These data provide a useful sequence ensemble
for studying which backbone positions are tightly constrained, but they do not
yet rank binding.

\section{Where a physics-aware score could enter}

The cleanest initial use of a new interface score is beside ProteinMPNN, after
RFdiffusion has fixed the backbone:
\[
\begin{aligned}
  \text{RFdiffusion backbones}
  &\longrightarrow \text{ProteinMPNN sequence ensemble} \\
  &\longrightarrow \text{all-atom reconstruction and interface scoring} \\
  &\longrightarrow \text{diverse shortlist}.
\end{aligned}
\]
This preserves the interpretation of each score. ProteinMPNN likelihood asks
whether a sequence is compatible with a backbone; the physical score asks
whether the resulting interface has favorable chemistry and geometry. An
initial score could include clash penalties, buried nonpolar surface,
unsatisfied polar groups, hydrogen-bond geometry, electrostatic complementarity,
and penalties for exposed hydrophobes.

Three levels of integration should be distinguished:
\begin{enumerate}
  \item \textbf{Reranking:} generate sequences normally, reconstruct side
  chains, and rank with the physical score. This is simplest and most auditable.
  \item \textbf{Biased sampling:} combine ProteinMPNN log probability with a
  differentiable or locally updated interface score during sequence sampling.
  This can improve efficiency but requires calibration of two unlike scores.
  \item \textbf{Joint retraining:} train or fine-tune a sequence model with the
  physical objective. This is the most ambitious option and needs reliable
  labels, negative examples, and controls against score exploitation.
\end{enumerate}

For a starter study, reranking is preferable because it lets us test whether
the physical score adds information beyond ProteinMPNN likelihood without
changing either pretrained model.

\section{Planned Week 35 deliverables}

\begin{enumerate}
  \item Trace one TYK2--4GIH forward pass through the released Boltz-2,
  Nesso-1, and FlashBind implementations and identify the exact information
  consumed by each affinity head.
  \item Add one architecture-matched perturbation for each affinity model, as
  defined in the Week 34 report.
  \item Present the RFdiffusion target, contig, hotspot tensor, noisy binder
  frames, and final backbone as an explicit input--output diagram.
  \item Present the ProteinMPNN residue graph, design mask, logits, sampled
  sequences, and position-wise entropy for the fixed insulin-receptor binder
  backbone.
  \item Define a small, prespecified physics-aware reranking score and test
  whether it selects information different from ProteinMPNN likelihood. No
  claim of affinity improvement will be made without an independent structural
  or experimental test.
\end{enumerate}

\section{Evidence boundary}

The insulin-receptor calculations currently establish execution,
conditioning behavior, stochastic structural diversity, and sequence diversity.
They do not establish that any generated sequence folds, binds the receptor,
modulates signaling, or has measurable affinity. The first useful scientific
advance is therefore not simply generating more candidates; it is defining an
independent and falsifiable selection procedure.

\bibliographystyle{unsrtnat}
\bibliography{references}
\end{document}
